\documentclass[11pt]{amsart}

\usepackage{amsmath,amssymb,amsthm}
\usepackage[margin=1in]{geometry}
\usepackage{hyperref}
\usepackage{mathtools}

\theoremstyle{plain}
\newtheorem{theorem}{Theorem}[section]
\newtheorem{proposition}[theorem]{Proposition}
\newtheorem{lemma}[theorem]{Lemma}
\newtheorem{corollary}[theorem]{Corollary}

\theoremstyle{definition}
\newtheorem{definition}[theorem]{Definition}
\newtheorem{example}[theorem]{Example}
\newtheorem{remark}[theorem]{Remark}

\newcommand{\Poly}{\mathbf{Poly}}
\newcommand{\Cat}{\mathbf{Cat}}
\newcommand{\Cathash}{\mathbf{Cat}^{\sharp}}
\newcommand{\Set}{\mathbf{Set}}
\newcommand{\Comod}{\mathrm{Comod}}
\newcommand{\PRA}{\mathrm{PRA}}
\newcommand{\CLP}{\mathrm{CLP}}
\newcommand{\Fam}{\mathrm{Fam}}
\renewcommand{\lhd}{\triangleleft}
\newcommand{\Store}[1]{\Delta #1}
\newcommand{\op}{^{\mathrm{op}}}
\newcommand{\yon}{\mathsf{y}}
\DeclareMathOperator{\const}{const}
\DeclareMathOperator{\ev}{ev}
\DeclareMathOperator{\im}{im}
\DeclareMathOperator{\Hom}{Hom}

\title[Polynomial bicomodules are parametric right adjoints]
{Polynomial Bicomodules are Parametric Right Adjoints\\
(after Spivak--Garner--Fairbanks),\\ with the worker/harness weld as a worked instance}

\author{MacBeth}
\address{Kodamai}
\date{18 October 2026}

\begin{document}

\begin{abstract}
Spivak, Garner and Fairbanks proved that, for polynomial comonoids (equivalently small
categories) $c$ and $d$, the category of $(c,d)$-bicomodules in $\Poly$ is equivalent to
the category of parametric right adjoint functors $c\text{-}\Set\to d\text{-}\Set$. We
restate this correspondence precisely, with its slice (aggregation) factorization, and
then run it on a single concrete bicomodule drawn from the semantics of stateful agents:
the worker/harness \emph{weld} $B=\Store{S}\lhd\Store{T}$ of two store comonads. We prove
that the weld is the \emph{regular} $(\Store{T},\Store{S})$-bicomodule and that it
represents the prafunctor $\bar F_B(V)=(T\times V)^S=T^S\cdot V^S$; its slice
factorization keeps the power $(-)^S$ as connected-limit content and aggregates over the
branch-map set $T^S$ by a coproduct $\Sigma$. We then locate the \emph{worker}
$A=\Store{(S\times T)}$ inside this picture: read as a $(\Store{T},\Store{S})$-bicomodule
of the same variance as $B$, it is exactly the constant-branch sub-bicomodule
$B'\subseteq B$, so worker and weld \emph{share} the content power $(-)^S$ and differ only
in their $\Sigma$-index, $T$ versus $T^S$. This decides a former conjecture of ours
(``worker $=$ parametric right adjoint''): affirmatively but with a sharp correction --- the
worker is not the whole weld-prafunctor but its constant-branch sub-bicomodule, with proper
inclusion exactly when $|S|\ge 2$ and $|T|\ge 2$.
\end{abstract}

\maketitle

\footnotetext{Supersedes the 2026-09-20 draft; corrects attribution and resolves the
former Conjecture~5.1.}

\section{Introduction}\label{sec:intro}

This note has two jobs. The first is expository: to state, carefully and with correct
attribution, the theorem of Spivak, Garner and Fairbanks \cite{SGF} identifying polynomial
bicomodules with parametric right adjoints. The second is a worked instance: to run that
theorem on one concrete bicomodule arising in the semantics of stateful, interacting
agents, and in doing so to settle a question we had earlier been able only to conjecture.

\subsection*{The correspondence}
In the polynomial-functor framework \cite{NiuSpivak}, a small category is the same datum as
a polynomial comonoid in $(\Poly,\lhd,\yon)$ \cite{AhmanUustalu}. The ``horizontal'' maps
between two such categories $c,d$ are the two-sided \emph{bicomodules}, and Spivak, Garner
and Fairbanks \cite{SGF} identify them concretely: a $(c,d)$-bicomodule is exactly a
\emph{parametric right adjoint} functor $c\text{-}\Set\to d\text{-}\Set$ between the
associated copresheaf categories. This is their Proposition~3.9(5) together with
Theorem~4.28. Moreover every prafunctor factors (Definition~3.5, equation~(21)) through a
slice, the first factor carrying the genuine right-adjoint (connected-limit) content and a
coproduct $\Sigma$ aggregating a tail that is not itself bicomodule data.

\subsection*{A note on attribution}
The result is sometimes loosely attributed to ``Garner 2019.'' That attribution is not
citable: it refers to a HoTTEST \emph{talk} (R.~Garner, 11 December 2019), an announcement,
not a written proof. The definitive, peer-reviewed written source is Spivak, Garner and
Fairbanks, \emph{Functorial Aggregation}, J.\ Pure Appl.\ Algebra \textbf{229}(2) (2025),
art.~107883; arXiv:2111.10968 \cite{SGF}. Garner is credited there, as a co-author of the
actual paper. We cite \cite{SGF} throughout and do not treat the 2019 talk as a proof.

\subsection*{The worked instance: the weld}
In models of stateful agents one carries a piece of local state $S$ by the \emph{store
comonad} $\Store{S}=S\cdot\yon^S$. Two agents with states $S$ and $T$ can be welded by
functor composition, $B=\Store{S}\lhd\Store{T}$ --- the ``store of stores'' used by
Braithwaite--Hedges--Mihejevs \cite{workersProof}. A competing discipline entangles the
states by the Dirichlet product into a single store,
$A=\Store{S}\otimes\Store{T}=\Store{(S\times T)}$, which we call the \emph{worker}. We
earlier conjectured (in the superseded 2026-09-20 draft) that the worker ``is'' the
parametric right adjoint content of the weld, with the weld's extra data as the
$\Sigma$-aggregation tail. We can now say precisely what is true.

The weld $B$ \emph{is} a bicomodule --- the regular $(\Store{T},\Store{S})$-bicomodule
(\S\ref{sec:weld}) --- and it represents the prafunctor $\bar F_B(V)=(T\times V)^S$, whose
factorization keeps the power $(-)^S$ as content and aggregates over the branch-map set
$T^S$. The worker $A$, read consistently as a bicomodule of the same variance, is the
\emph{constant-branch sub-bicomodule} $B'\subseteq B$ (\S\ref{sec:worker}). Worker and weld
share the content power $(-)^S$; they differ only in their $\Sigma$-index, $T$ versus $T^S$.
So the former conjecture is settled affirmatively in a precise and corrected form: the
worker is a bicomodule, and it is not the whole weld-prafunctor but its constant-branch
sub-bicomodule.

\subsection*{Status of the claims}
The correspondence of \S\ref{sec:corr} is \cite{SGF}'s established theorem; we state it and
sketch its proof in our own words, and claim no part of it. The weld facts of
\S\S\ref{sec:weld}--\ref{sec:worker} are our own results; the regular-bicomodule structure
and the represented prafunctor $\bar F_B(V)=(T\times V)^S$ have additionally been
peer-reviewed by a referee (Rick), and the worker-as-constant-branch-slice comparison is
proved and machine-checked at small shapes. Trust grades are flagged at each statement.

\subsection*{Outline}
Section~\ref{sec:prelim} fixes polynomial conventions, the store comonad, parametric right
adjoints, and the bicomodule definition. Section~\ref{sec:corr} states the SGF
correspondence, its slice factorization, and the handedness convention.
Section~\ref{sec:weld} proves the weld is the regular bicomodule and computes its
prafunctor. Section~\ref{sec:worker} locates the worker as the constant-branch
sub-bicomodule and decides the former conjecture. Section~\ref{sec:conc} draws the
consequence for the compositional-correctness program.

\section{Preliminaries}\label{sec:prelim}

We work in the category $\Poly$ of polynomial functors $\Set\to\Set$, following the
conventions of \cite{NiuSpivak}.

\begin{definition}[polynomials and their maps]\label{def:poly}
A \emph{polynomial} is a coproduct of representables $p=\sum_{s\in S_p}\yon^{p[s]}$, acting
on sets by $p(X)=\sum_{s\in S_p}X^{p[s]}$; we call $S_p$ its \emph{positions} (or shapes)
and $p[s]$ its set of \emph{directions} at $s$. A morphism $f\colon p\to q$ is a pair
$f=(f_1,f^{\sharp})$ with $f_1\colon S_p\to S_q$ forward on positions and
$f^{\sharp}_s\colon q[f_1 s]\to p[s]$ \emph{backward} on directions. Composition is
$(g\circ f)_1=g_1 f_1$ and $(g\circ f)^{\sharp}_s=f^{\sharp}_s\circ g^{\sharp}_{f_1 s}$:
backward maps compose in reverse. Under the full embedding $\Poly\hookrightarrow[\Set,\Set]$
a morphism of polynomials \emph{is} a natural transformation.
\end{definition}

$\Poly$ carries two monoidal structures we need.

\begin{definition}[the two products]\label{def:products}
The \emph{composition product} is $(p\lhd q)(X)=p(q(X))$, with unit $\yon$; it is strictly
associative. The \emph{Dirichlet product} $p\otimes q$ has positions $S_p\times S_q$ and
directions $p[s]\times q[t]$ at $(s,t)$, with unit $\yon$.
\end{definition}

\begin{definition}[polynomial comonoids are categories]\label{def:comonoid}
A comonoid for $(\Poly,\lhd,\yon)$ is the same data as a small category
\cite{AhmanUustalu,NiuSpivak}: positions are objects, directions at $s$ are the morphisms
out of $s$, the counit $\varepsilon$ picks out identities and the comultiplication $\delta$
composes. The framed bicategory of polynomial comonoids is written $\Cathash$
\cite{NiuSpivak}. A polynomial comonoid $c$ presents a category we also call $c$, with
copresheaf category $c\text{-}\Set$.
\end{definition}

\begin{definition}[store comonad]\label{def:store}
For a set $S$, the \emph{store} (or costate) polynomial is $\Store{S}:=S\cdot\yon^{S}$ ---
positions $S$, with $\Store{S}[s]=S$ at every position --- carrying the store comonad
$\Store{S}(X)=S\times X^{S}$. Its comultiplication is $\delta_S(s,\varphi)=(s,\,i\mapsto
(i,\varphi))$ and counit $\varepsilon_S(s,\varphi)=\varphi(s)$. As a category $\Store{S}$ is
the \emph{indiscrete} (codiscrete) category on $S$: exactly one morphism between any two
objects. For $S\ne\emptyset$ it is equivalent to the terminal category $1$, so
$\Store{S}\text{-}\Set\simeq\Set$.
\end{definition}

The two objects of study on the SGF side are parametric right adjoints and bicomodules.

\begin{definition}[parametric right adjoint]\label{def:pra}
Let $C$ have a terminal object $1$. A functor $F\colon C\to D$ is a \emph{parametric right
adjoint} (PRA), or \emph{prafunctor}, if the first factor of the canonical slice
factorization
\[
C \xrightarrow{\;F/1\;} D/F1 \xrightarrow{\;U\;} D
\]
is a right adjoint, where $(F/1)(c)=(Fc\to F1)$ and $U$ forgets the slice down to $D$
\cite[Def.~3.5, eq.~(21)]{SGF}. For endofunctors of $\Set$ this is equivalent to preserving
connected limits, equivalently to being polynomial \cite[Prop.~2.5]{SGF}.
\end{definition}

\begin{definition}[bicomodule]\label{def:bicomod}
Let $c,d$ be polynomial comonoids. A \emph{$(c,d)$-bicomodule} \cite[Def.~4.2]{SGF} is a
polynomial $m$ with a left $d$-coaction $\lambda\colon m\to d\lhd m$ and a right
$c$-coaction $\rho\colon m\to m\lhd c$,
\[
d\lhd m \xleftarrow{\;\lambda\;} m \xrightarrow{\;\rho\;} m\lhd c,
\]
satisfying the left comodule laws (counit and coassociativity, \cite[(34)]{SGF}), the right
comodule laws (\cite[(35)]{SGF}), and the compatibility law
$(\lambda\lhd c)\circ\rho=(d\lhd\rho)\circ\lambda$ (\cite[(36)]{SGF}). Bicomodules between
$c$ and $d$ are the ``horizontal'' (loose) maps of $\Cathash$.
\end{definition}

\section{The Spivak--Garner--Fairbanks correspondence}\label{sec:corr}

We restate the results of \cite{SGF} this note leans on, in our own words. All are
established; we claim no part of them.

\begin{theorem}[{\cite[Thm.~4.28]{SGF}}]\label{thm:sgf-main}
The framed bicategories $\Cathash$ and $\Comod(\Poly)$ are equivalent. Under the
equivalence: objects are small categories, i.e.\ polynomial comonoids; the horizontal
(loose) maps from $c$ to $d$ are the $(c,d)$-bicomodules; the vertical (tight) maps are the
comonoid homomorphisms, i.e.\ retrofunctors (\emph{not} ordinary functors); and the
$2$-cells are bicomodule morphisms.
\end{theorem}

The pointwise identification of the horizontal maps is a six-way equivalence, of which the
piece we use is item~(5): parametric right adjoints.

\begin{proposition}[{\cite[Prop.~3.9(5), with Thm.~4.28]{SGF}}]\label{prop:sixway}
For polynomial comonoids $c,d$ there is an equivalence of categories
\[
\Comod(\Poly)(c,d)\;\simeq\;\PRA\bigl(c\text{-}\Set,\,d\text{-}\Set\bigr)
\;=\;\CLP\bigl(c\text{-}\Set,\,d\text{-}\Set\bigr),
\]
where $\CLP$ denotes the connected-limit-preserving functors. That is, a
$(c,d)$-bicomodule \emph{is} a parametric right adjoint functor $c\text{-}\Set\to
d\text{-}\Set$. (Prop.~3.9 lists six equivalent descriptions of the hom-category, with the
bicomodule description supplied via Theorem~\ref{thm:sgf-main}; items also include functors
$d\to\Set[c]$ and $d\to\Fam((c\text{-}\Set)\op)$.)
\end{proposition}

\begin{proof}[Proof sketch, following \cite{SGF}]
A $(c,d)$-bicomodule $m$ determines a functor $c\text{-}\Set\to d\text{-}\Set$ by right Kan
extension / coend against its carrier, using the two coactions to make the output a
$d$-copresheaf and to integrate over $c$. That this functor is a parametric right adjoint is
the content of \cite[\S2--\S3]{SGF}: polynomial endofunctors of $\Set$ are exactly the
connected-limit-preserving ones (\cite[Prop.~2.5]{SGF}), and the dependent-sum indexing of a
bicomodule carrier realises the slice factorization of Definition~\ref{def:pra}. Conversely a
prafunctor is reconstructed from the slice data $F/1$ and the indexing set $F1$, which assemble
into a bicomodule carrier with its coactions; composition is covariant and matches ordinary
functor composition via the coreflexive-equalizer composite $\odot_d$, $\Poly$ preserving
coreflexive equalizers (\cite[Lem.~2.19]{SGF}). The two passages are mutually inverse up to
coherent isomorphism.
\end{proof}

The factorization we actually compare against is the slice factorization of a single
prafunctor, because it is what performs the ``content versus aggregation'' split.

\begin{proposition}[slice factorization, {\cite[Def.~3.5, eq.~(21)]{SGF}}]\label{prop:factor}
Every prafunctor $F\colon c\text{-}\Set\to\Set$ factors as
\[
c\text{-}\Set \xrightarrow{\;F/1\;} \Set/F1 \xrightarrow{\;\Sigma\;} \Set,
\]
where $F/1$ is a \emph{right adjoint} --- the genuine polynomial/bicomodule content --- and
$\Sigma$ is the coproduct over the index set $F1$, a \emph{left adjoint} ``aggregation
tail'' that carries no bicomodule data. (The factorization is the datum named in
Definition~3.5, equation~(21); it is elaborated inside the proof of \cite[Prop.~3.6]{SGF},
which is itself the single-variable equivalence statement.)
\end{proposition}

The moral is a content/aggregation split: the bicomodule content lives in the right-adjoint
factor $F/1$, and the coproduct $\Sigma$ is a tail the correspondence does not see as
bicomodule data. A second, equivalent presentation of a prafunctor is the data-migration
form, which we record because it names the pieces the weld will exhibit.

\begin{proposition}[data-migration form, {\cite[Prop.~3.20]{SGF}}]\label{prop:datamig}
Every prafunctor factors as $\Sigma_T\circ\Pi_\pi\circ\Delta_S$: a reindexing
$\Delta_S$ (pullback), a dependent product $\Pi_\pi$ (the connected-limit content), and a
dependent sum $\Sigma_T$ (the aggregation tail).
\end{proposition}

\begin{remark}[handedness]\label{rem:handedness}
We use \cite{SGF}'s convention throughout: a $(c,d)$-bicomodule $m\colon c\to d$
corresponds to a parametric right adjoint $c\text{-}\Set\to d\text{-}\Set$
(Proposition~\ref{prop:sixway}). The reader should be warned that Lynch, Shapiro and Spivak,
\emph{All concepts are $\Cathash$} \cite{LSS}, adopt the \emph{reversed} handedness, sending a
bicomodule to a functor $d\text{-}\Set\to c\text{-}\Set$. It is the same theorem under the
opposite variance on bicomodules; any concrete identification must pin the convention first.
We pin \cite{SGF}'s.
\end{remark}

\section{The weld is the regular bicomodule}\label{sec:weld}

We now run the correspondence on the concrete bicomodule of interest. Fix nonempty sets
$S,T$ and the store comonoids $\Store{S},\Store{T}$ of Definition~\ref{def:store}. The
\emph{weld} is the composition product
\[
B\;:=\;\Store{S}\lhd\Store{T},
\]
with positions $\{(s,g):s\in S,\ g\colon S\to T\}=S\times T^{S}$ and directions $S\times T$
everywhere, so a position of $B$ carries a whole \emph{branch map} $g\colon S\to T$.

\begin{lemma}[regular bicomodule]\label{lem:regular}
In any monoidal category $(\mathcal V,\lhd,\yon)$, for comonoids
$c=(c,\varepsilon_c,\delta_c)$ and $d=(d,\varepsilon_d,\delta_d)$, the object $d\lhd c$ is a
$(c,d)$-bicomodule with
\[
\lambda:=\delta_d\lhd c\colon d\lhd c\to d\lhd d\lhd c=d\lhd(d\lhd c),
\qquad
\rho:=d\lhd\delta_c\colon d\lhd c\to d\lhd c\lhd c=(d\lhd c)\lhd c.
\]
\end{lemma}

\begin{proof}
The left comodule laws \cite[(34)]{SGF} are the comonoid laws of $d$ ---
$(\varepsilon_d\lhd d)\circ\delta_d=\mathrm{id}$ and
$(\delta_d\lhd d)\circ\delta_d=(d\lhd\delta_d)\circ\delta_d$ --- whiskered on the right by
$\lhd c$; the right laws \cite[(35)]{SGF} are the comonoid laws of $c$ whiskered on the left
by $d\lhd$. For compatibility \cite[(36)]{SGF}, $\lambda$ acts on the first ($d$) factor and
$\rho$ on the last ($c$) factor; since $\lhd$ is a bifunctor, the two whiskerings commute by
the interchange law, and both $(\lambda\lhd c)\circ\rho$ and $(d\lhd\rho)\circ\lambda$ equal
the single map $\delta_d\lhd\delta_c\colon d\lhd c\to d\lhd d\lhd c\lhd c$.
\end{proof}

\begin{theorem}[weld as regular bicomodule]\label{thm:weld-regular}
The weld $B=\Store{S}\lhd\Store{T}$ is the regular $(\Store{T},\Store{S})$-bicomodule: take
$c=\Store{T}$, $d=\Store{S}$, carrier $d\lhd c=\Store{S}\lhd\Store{T}=B$, with
$\lambda=\delta_{\Store{S}}\lhd\Store{T}$ and $\rho=\Store{S}\lhd\delta_{\Store{T}}$. Hence
$B$ is a horizontal map of $\Cathash$, i.e.\ a prafunctor $\Store{T}\text{-}\Set\to
\Store{S}\text{-}\Set$.
\end{theorem}

\begin{proof}
Immediate from Lemma~\ref{lem:regular} with $c=\Store{T}$, $d=\Store{S}$, since
$\Store{S},\Store{T}$ are comonoids (Definition~\ref{def:store}). The pair $(c,d)$ in the
sense of Proposition~\ref{prop:sixway} is $(\Store{T},\Store{S})$.
\end{proof}

\begin{remark}[trust]\label{rem:trust-weld}
Theorem~\ref{thm:weld-regular} is \emph{proved} and \emph{peer-reviewed} (referee: Rick,
clause~(1)/Thm~2.2). All five bicomodule axioms were additionally machine-checked as
equalities of natural transformations for
$(|S|,|T|)\in\{(2,2),(2,3),(3,2),(1,2),(2,1)\}$. Lemma~\ref{lem:regular} is standard
category theory (a $\lhd$-product of comonoids is the regular bicomodule between them); the
content here is its use to pin the pair $(\Store{T},\Store{S})$ for the weld, with no
ad hoc choice required.
\end{remark}

The theorem removes the only genuine type-level worry about the weld: although
$\Store{S}\lhd\Store{T}$ is a priori a \emph{product} of two comonoids rather than a map
\emph{between} two fixed ones, the regular-bicomodule construction turns the $\lhd$-product
of any two comonoids into a bicomodule between them.

\subsection*{The represented prafunctor}
Because $\Store{S}$ is the indiscrete category on $S$, it is equivalent to the terminal
category, so $\Store{S}\text{-}\Set\simeq\Set$ (Definition~\ref{def:store}); likewise for
$\Store{T}$. Right-adjointness and the slice factorization are equivalence-invariant, so we
transport the weld prafunctor to an endofunctor $\bar F_B$ of $\Set$.

\begin{theorem}[represented prafunctor]\label{thm:weld-pra}
$\displaystyle \bar F_B(V)\;\cong\;(T\times V)^{S}\;=\;T^{S}\times V^{S}
\;=\;\sum_{g\in T^{S}}V^{S}.$
As a polynomial endofunctor of $\Set$, $\bar F_B=T^{S}\cdot\yon^{S}$: positions $T^{S}$,
directions $S$. Concretely $\bar F_B(X)=\{(g,h):g\colon S\to T,\ h\colon S\to X\}$.
\end{theorem}

\begin{proof}
The regular bicomodule $d\lhd c$ is the composite, over the unit $\yon$, of two one-sided
comodules,
\[
\Store{T}\xrightarrow{\;G\;}\yon\xrightarrow{\;H\;}\Store{S},
\qquad\text{carrier } H\lhd G=\Store{S}\lhd\Store{T}=B,
\]
where $G$ is the right $\Store{T}$-comodule with carrier $\Store{T}$ and $\rho=\delta_{\Store{T}}$,
and $H$ is the left $\Store{S}$-comodule with carrier $\Store{S}$ and
$\lambda=\delta_{\Store{S}}$. Bicomodule composition over the unit $\yon$ is just $\lhd$
(no equalizer is needed), so prafunctors compose and $\bar F_B=\bar F_H\circ\bar F_G$.

\emph{The factor $\bar F_G\colon\Store{T}\text{-}\Set\to\Set$.} As a $\Store{T}$-variable
polynomial (the $\Set[c]$ description of \cite[Prop.~3.9]{SGF}) the right comodule with
carrier $\Store{T}$ is $F_G(X)=\sum_{t\in T}\Store{T}\text{-}\Set(\yon_t,X)=\sum_{t\in
T}X(t)$, where $\yon_t$ is the representable at $t$. In the indiscrete category $\yon_t$ is
the terminal copresheaf, so $\Store{T}\text{-}\Set(\yon_t,X)=\lim X$; hence
$F_G(X)=\sum_{t\in T}\lim X$, and after transport $\bar F_G(V)=T\times V$.

\emph{The factor $\bar F_H\colon\Set\to\Store{S}\text{-}\Set$.} A left $\Store{S}$-comodule
is a functor $\Store{S}\to\Poly$ (the Prop.~5.7 dictionary of \cite{SGF}); the regular one
(carrier $\Store{S}$, $\lambda=\delta_{\Store{S}}$) assigns to each object $a\in S$ the
component polynomial $\yon^{S}$ (read off $\delta_S$: one position over $a$, with $S$
directions). So $F_H(W)(a)=(\yon^{S})(W)=W^{S}$, and after transport $\bar F_H(W)=W^{S}$.

Therefore $\bar F_B(V)=\bar F_H(\bar F_G(V))=(T\times V)^{S}=T^{S}\times V^{S}$, and
$(T\times V)^{S}=\{S\to T\times V\}=\{(g,h):g\colon S\to T,\ h\colon S\to V\}$.
\end{proof}

\begin{remark}[trust, and the point of the computation]\label{rem:trust-pra}
Theorem~\ref{thm:weld-pra} is \emph{proved} and \emph{peer-reviewed} (Rick re-derived
$\bar F_B$ on seven shapes). Note that $\bar F_B=T^{S}\yon^{S}$ is \emph{neither} the
underlying polynomial of the worker $A=(S\times T)\yon^{S\times T}$ \emph{nor} that of the
carrier $B=(S\cdot T^{S})\yon^{S\times T}$: the prafunctor represented by a bicomodule is not
the underlying polynomial functor of its carrier. (Cardinalities at $|S|=|T|=2$, $V=2$:
$\bar F_B=16$, $A$-underlying $=64$, $B$-underlying $=128$.)
\end{remark}

Applying Proposition~\ref{prop:factor} to $\bar F_B(V)=\sum_{g\in T^{S}}V^{S}$:
\[
\bar F_B(1)=T^{S},\qquad
(\bar F_B/1)(V)=\bigl(V^{S}\bigr)_{g\in T^{S}}=\Delta_{T^{S}}\circ(-)^{S},\qquad
\Sigma\colon \Set/T^{S}\to\Set.
\]
The content factor $\bar F_B/1$ is a right adjoint --- the power $(-)^{S}$ (right adjoint to
the copower $S\times(-)$) followed by the diagonal $\Delta_{T^{S}}$ (right adjoint to
$\Sigma$). The connected-limit content is the \textbf{power $(-)^{S}$}; the aggregation tail
is the coproduct $\Sigma$ over the \textbf{branch-map set $T^{S}=\Hom(S,T)$}.

\section{The worker is the constant-branch sub-bicomodule}\label{sec:worker}

The \emph{worker} is the Dirichlet entanglement
\[
A\;:=\;\Store{S}\otimes\Store{T}\;=\;\Store{(S\times T)},
\]
with positions $S\times T$ and directions $S\times T$ everywhere. A priori $A$ is a
\emph{$0$-cell} (a comonoid/category), whereas the weld content is a \emph{$1$-cell}
($\Store{T}\to\Store{S}$); to compare them we must read $A$ as a
$(\Store{T},\Store{S})$-bicomodule of the \emph{same variance} as $B$. There is a canonical
such reading, and it sits visibly inside $B$.

\begin{definition}[constant-branch slice]\label{def:Bprime}
Writing a position of $B$ as $(s,g)$ with $s\in S$, $g\in T^{S}$, let
\[
B'\;:=\;\{(s,g)\in B:\ g\in\const(T)\}, \qquad \const\colon T\hookrightarrow T^{S},\ t\mapsto\const_t,
\]
with directions inherited ($S\times T$ per position). As a polynomial, $B'$ has positions
$S\times\const(T)\cong S\times T$ and directions $S\times T$, i.e.\
$B'=(S\times T)\cdot\yon^{S\times T}=\Store{(S\times T)}=A$ on the nose.
\end{definition}

\begin{lemma}[$B'$ is a sub-bicomodule]\label{lem:subbicomod}
$B'$ is closed under the weld coactions $\lambda=\delta_{\Store{S}}\lhd\Store{T}$ and
$\rho=\Store{S}\lhd\delta_{\Store{T}}$; hence $(B',\lambda|,\rho|)$ is a
$(\Store{T},\Store{S})$-sub-bicomodule of $B$.
\end{lemma}

\begin{proof}
Both coactions leave the branch-map coordinate $g=(t_s)_{s\in S}$ \emph{pointwise}
unchanged, only reorganising positions. For $\rho=\Store{S}\lhd\delta_{\Store{T}}$: the
inner $\delta_T(t,\psi)=(t,\,j\mapsto(j,\psi))$ keeps the first coordinate $t$, so the
$B$-layer of $\rho(s,g)$ has the same $g$; a constant $g$ stays constant. For
$\lambda=\delta_{\Store{S}}$: $\delta_S(s,\varphi)=(s,\,i\mapsto(i,\varphi))$ duplicates the
whole $\Store{T}$-ground with the same $\varphi$, hence the same $g$; constant stays
constant. The counit and (co)associativity laws restrict because they are the comonoid laws
of $\Store{S},\Store{T}$ whiskered, and $B'$ is full on positions with inherited directions.
Machine-checked ($\rho$-closed, $\lambda$-closed) at
$(|S|,|T|)\in\{(2,2),(2,3),(3,2),(1,2),(2,1)\}$.
\end{proof}

\begin{proposition}[the worker's canonical reading]\label{prop:sigma}
The map $\sigma\colon A\to B$ with $\sigma_1(s,t)=(s,\const_t)$ on positions and
$\sigma^{\sharp}=\mathrm{id}_{S\times T}$ on directions is an isomorphism of
$(\Store{T},\Store{S})$-bicomodules $A\cong B'$. Equivalently, $A$ carries the
$(\Store{T},\Store{S})$-bicomodule structure obtained by cobase change of the regular
$(A,A)$-bicomodule along the comonoid projections $\pi_S\colon A\to\Store{S}$ (left) and
$\pi_T\colon A\to\Store{T}$ (right) --- the indiscrete functors induced by the set
projections --- and this structure coincides with $B'$.
\end{proposition}

\begin{proof}
$\sigma_1$ is the bijection $S\times T\cong S\times\const(T)=B'$-positions, and
$\sigma^{\sharp}=\mathrm{id}_{S\times T}$ matches directions. By Lemma~\ref{lem:subbicomod}
the coactions act on $B'$ exactly as they do on $A$ under this identification (both leave the
$g$-value alone), so $\sigma$ commutes with $\lambda,\rho$ and is a bicomodule isomorphism
onto its image $B'$. The sub-bicomodule description is primary and fully verified; the
cobase-change description is the conceptual justification that $B'$ is the \emph{canonical}
reading (not an arbitrary sub-object), and agrees with it on carriers and positions.
\end{proof}

\begin{theorem}[worker versus weld: a $\Sigma$-index separation]\label{thm:separation}
Read $A$ and $B$ as $(\Store{T},\Store{S})$-bicomodules via Proposition~\ref{prop:sigma}
($A\cong B'\subseteq B$). Then
\[
\bar F_B(V)=\sum_{g\in T^{S}}V^{S}=T^{S}\cdot V^{S},
\qquad
\bar F_A(V)=\sum_{g\in\const(T)}V^{S}=T\cdot V^{S}.
\]
Worker and weld have the \emph{same} content power $(-)^{S}$ (directions $S$); they differ
\emph{only} in the $\Sigma$-aggregation index, $T^{S}$ (all branch maps) for the weld versus
$T\cong\const(T)$ (constant branch maps) for the worker. Hence $\bar F_A$ is a
sub-prafunctor of $\bar F_B$, with $\bar F_A\cong\bar F_B$ \emph{iff} $|S|=1$ or $|T|=1$, and
a \emph{proper} inclusion exactly when $|S|\ge 2$ and $|T|\ge 2$.
\end{theorem}

\begin{proof}
$\bar F_B=\sum_{g\in T^{S}}V^{S}$ is Theorem~\ref{thm:weld-pra}; its $\Sigma$-index is the set
of $B$-positions modulo the left $\Store{S}$-collapse, namely $T^{S}$ (the $S$-head is killed
by indiscreteness), and the content exponent is the surviving direction $S$ (the
$T$-direction is killed by the right $\Store{T}$-collapse). By Proposition~\ref{prop:sigma},
$A\cong B'$ is $B$ restricted on positions to $\const(T)\subseteq T^{S}$ with directions
unchanged; the same recipe restricts the $\Sigma$-index to $\const(T)$ and leaves the content
power $(-)^{S}$ untouched, giving $\bar F_A(V)=\sum_{t\in T}V^{S}=T\cdot V^{S}$. Equality of
the two functors for all $V$ is equality of polynomials $T^{S}\yon^{S}$ and $T\yon^{S}$, i.e.\
$|T^{S}|=|T|$, i.e.\ $m^{n}=m$ (with $n=|S|$, $m=|T|\ge 1$), i.e.\ $n=1$ or $m=1$. Otherwise
$T\subsetneq T^{S}$ and the inclusion $\bar F_A\hookrightarrow\bar F_B$ is proper.
\end{proof}

\begin{remark}[what this corrects]\label{rem:correct}
The content \emph{exponents agree} --- both are $(-)^{S}$. An earlier, superseded argument
read $A$ \emph{without} the right $\Store{T}$-coaction and so reported content exponent $S$
for the weld against $S\times T$ for the worker, concluding a mismatch already at $|S|=1$. But
at $|S|=1$ one has $T^{S}=T$ and $B=\yon\lhd\Store{T}=\Store{T}=\Store{(1\times T)}=A$ on the
nose: worker $=$ weld there, so the mismatch claim was an artefact of the inconsistent reading.
With the consistent $(\Store{T},\Store{S})$-reading the $T$-directions of $A$ collapse into the
exponent $S$, and the true separation is in the $\Sigma$-index, $T$ versus $T^{S}$, nonzero
exactly where non-constant branch maps exist: $|S|\ge 2$ and $|T|\ge 2$.
\end{remark}

\subsection*{What the retract forgets}
The states-entanglement discipline comes with a retraction $r\colon B\to A$,
$r_1(s,g)=(s,g(s))$ (self-evaluation), splitting $\sigma$; the idempotent $e=\sigma\circ r$
has $e(s,g)=(s,\const_{g(s)})$.

\begin{theorem}[the retract's keep/discard]\label{thm:retract}
The idempotent $e\colon S\times T^{S}\to S\times T^{S}$ has image $S\times\const(T)\cong
S\times T$ (the worker positions). For each current input $s$, $e$ keeps the value
$g(s)\in T$ and discards the restriction $g|_{S\setminus\{s\}}\in T^{S\setminus\{s\}}$, a
fibre of size $|T|^{|S|-1}$. Concretely the fibre of $e$ over $(s,\const_t)$ is
$\{g:g(s)=t\}\cong T^{S\setminus\{s\}}$.
\end{theorem}

\begin{proof}
$e\circ e=e$ since $\const_{g(s)}$ evaluated at $s$ is $g(s)$; the image is the constant
slice. Two maps $g,g'$ satisfy $e(s,g)=e(s,g')$ iff $g(s)=g'(s)$; fixing $g(s)=t$ leaves
$S\setminus\{s\}$ free, so the fibre is $T^{S\setminus\{s\}}$. Idempotency, image, and all
fibre sizes machine-checked at the five shapes above.
\end{proof}

\begin{remark}[the $\Sigma$-tail, stated honestly]\label{rem:tail}
The weld's $\Sigma$-aggregation index is the \emph{full} branch-map set
$T^{S}=\const(T)\sqcup\mathrm{NC}$, with $|\const(T)|=m$ and $|\mathrm{NC}|=m^{n}-m$
(non-constant maps). It is \emph{not} correct to say ``$T^{S}=$ the non-constant maps''
(a cardinality contradiction: $m^{n}\ne m^{n}-m$), nor that the retract forgets exactly the
non-constant maps: per input $s$ the retract forgets the off-diagonal fibre
$T^{S\setminus\{s\}}$ of size $m^{n-1}$, which is neither $\mathrm{NC}$ ($m^{n}-m$) nor $T^{S}$
($m^{n}$). At $(|S|,|T|)=(2,3)$ these are $3$ versus $6$ versus $9$, all distinct. What is
correct: the weld's $\Sigma$-tail \emph{is} the full branch-map set $T^{S}$, and the retract
realises its projection onto the worker's constant slice $\const(T)\cong T$ by ``evaluate at
the current input, then make constant.''
\end{remark}

\begin{remark}[no reversal of adjointness]\label{rem:no-reversal}
The superseded draft worried that the retract discards ``exponential, hence
right-adjoint-flavoured'' data $T^{S}$ while the SGF factorization discards a coproduct.
There is no paradox. $T^{S}$ appears as the \emph{indexing set} of the left-adjoint coproduct
$\Sigma_{T^{S}}\dashv\Delta_{T^{S}}$, which is a left adjoint for any index set, exponential or
not. The genuinely right-adjoint (connected-limit) content $(-)^{S}$ is \emph{kept} in the
$F/1$ factor, exactly as \cite{SGF} predicts, and is shared by worker and weld.
\end{remark}

\subsection*{Resolution of the former Conjecture~5.1}
The superseded 2026-09-20 draft conjectured that ``a worker is a parametric right adjoint /
bicomodule, and the retract kernel is the aggregation tail.'' We can now decide it precisely.

\begin{theorem}[decision of the former conjecture]\label{thm:decision}
Let $S,T$ be nonempty.
\begin{enumerate}
\item[\textup{(1)}] \textbf{(existence --- true, peer-reviewed.)} The weld
$B=\Store{S}\lhd\Store{T}$ is a bicomodule: the regular $(\Store{T},\Store{S})$-bicomodule
(Theorem~\ref{thm:weld-regular}), representing the prafunctor $\bar F_B(V)=(T\times V)^{S}$
(Theorem~\ref{thm:weld-pra}).
\item[\textup{(2)}] \textbf{(identification --- corrected.)} The worker $A$, read as the
$(\Store{T},\Store{S})$-bicomodule $\sigma(A)=B'\subseteq B$, is \emph{not} the full
weld-prafunctor but its \emph{constant-branch sub-bicomodule}. Worker and weld share the
right-adjoint content power $(-)^{S}$ and differ only in the $\Sigma$-index ($T$ versus
$T^{S}$). The inclusion is proper exactly on $|S|\ge 2\wedge|T|\ge 2$ (witness $(2,2)$ at
$V=2$: $\bar F_A=8$, $\bar F_B=16$); worker $=$ weld on the nose when $|S|=1$ or $|T|=1$.
\item[\textup{(3)}] \textbf{(tail --- true.)} The weld's $\Sigma$-aggregation tail is the
coproduct over the branch-map set $T^{S}=\Hom(S,T)$; the retract projects this index onto the
worker's constant slice (Remark~\ref{rem:tail}).
\end{enumerate}
In particular the former conjecture holds in the corrected form: \emph{the worker is a
parametric right adjoint --- the constant-branch sub-bicomodule of the weld --- not the whole
weld-prafunctor.}
\end{theorem}

\begin{remark}[trust]\label{rem:trust-decision}
Clause~(1) (Theorems~\ref{thm:weld-regular}, \ref{thm:weld-pra}) is \emph{peer-reviewed}
(Rick). The worker-as-constant-branch-slice content of clauses~(2),~(3)
(Definition~\ref{def:Bprime}, Lemma~\ref{lem:subbicomod},
Propositions~\ref{prop:sigma}, Theorems~\ref{thm:separation}, \ref{thm:retract}) is
\emph{proved} and machine-checked at the five shapes
$(|S|,|T|)\in\{(2,2),(2,3),(3,2),(1,2),(2,1)\}$; the full coaction-by-coaction identification
of the two bicomodule structures of Proposition~\ref{prop:sigma} for general $S,T$ is not
written out here, as the comparison only uses $A\cong B'$.
\end{remark}

\section{What this buys the program}\label{sec:conc}

The weld now has a precise home in $\Poly$: it \emph{is} a data-migration prafunctor
$\Store{T}\text{-}\Set\to\Store{S}\text{-}\Set$, $X\mapsto(T\times X)^{S}$ --- ``read the
whole input state $S$; return, for each input, an output token in $T$ and a value in $X$,''
with the branch-map choice aggregated over $T^{S}$ by $\Sigma$. This places the \emph{state}
axis of the compositional-correctness program inside the very framed bicategory $\Cathash$
that already houses the directed (retrofunctor) axis.

But the identification ``worker $=$ parametric right adjoint content'' is \emph{false} on
$|S|,|T|\ge 2$: state-entanglement (the worker $A$) and data-aggregation (the weld's
$\Sigma$-tail) are genuinely different collapses of the same content power $(-)^{S}$. The
worker keeps a branch choice per output token $t\in T$; the weld keeps a branch choice per
\emph{function} $g\in T^{S}$. Keeping the two distinct is a semantic requirement of
agent-orchestration, not a bug --- and the exact locus where it matters is $|S|,|T|\ge 2$.
That is the honest, sharpened dividend of running the SGF correspondence on this one
bicomodule.

\begin{thebibliography}{9}

\bibitem{AhmanUustalu}
D.~Ahman and T.~Uustalu,
\emph{Directed containers as categories},
in \emph{Proc.\ MSFP 2016}, EPTCS 207, 2016, pp.~89--98.
arXiv:1604.01187.

\bibitem{LSS}
O.~Lynch, B.~T.~Shapiro and D.~I.~Spivak,
\emph{All concepts are $\mathbf{Cat}^{\sharp}$},
arXiv:2305.02571, 2023.

\bibitem{NiuSpivak}
N.~Niu and D.~I.~Spivak,
\emph{Polynomial Functors: A Mathematical Theory of Interaction},
arXiv:2312.00990, 2023.

\bibitem{SGF}
D.~I.~Spivak, R.~Garner and A.~Fairbanks,
\emph{Functorial aggregation},
J.\ Pure Appl.\ Algebra \textbf{229} (2025), no.~2, art.~107883.
arXiv:2111.10968.

\bibitem{workersProof}
MacBeth,
\emph{The worker/harness weld is the regular bicomodule} and
\emph{Worker vs.\ weld content: the $\Sigma$-index separation},
Kodamai internal proof notes, 2026-10-16 and 2026-10-18.
Prior art: J.~Braithwaite, J.~Hedges and Z.~Mihejevs (ACT 2026), graded
$\lhd$-composition of store comonads.

\end{thebibliography}

\end{document}
